Pairwise sequence alignment is a  very time-consuming step in common bioinformatics pipelines. Speeding up this step requires heuristics, efficient implementations, and/or hardware acceleration. A promising candidate for all of the above is the recently proposed GenASM algorithm. We identify and address three inefficiencies in the GenASM algorithm: it has a high amount of data movement, a large memory footprint, and does some unnecessary work.

We propose \defn\toolname, a fast and memory-frugal genomic sequence aligner. \toolname includes three novel algorithmic improvements which reduce the  data movement, memory footprint, and the number of operations in the GenASM algorithm. We provide efficient open-source implementations of the \toolname algorithm for CPUs and GPUs, which demonstrate the significant benefits of our algorithmic improvements.
For long reads, the CPU version of \toolname achieves a \kswcpuspeedup, \edlibcpuspeedup, and \cpuoptspeedup speedup over KSW2, Edlib, and a CPU implementation of GenASM, respectively. The GPU version of \toolname achieves a \cpugpuspeedup \kswgpuspeedup, \edlibgpuspeedup, \darwingpuspeedup and \gpuoptspeedup speedup over the CPU version of \toolname, KSW2, Edlib, Darwin-GPU, and a GPU implementation of GenASM, respectively. We estimate an ASIC implementation of \toolname to use 3.6\x less chip area and 2.1\x less power than a GenASM ASIC while maintaining the same throughput. Further, we systematically analyze the throughput and accuracy behavior of GenASM and \toolname under various configurations. As the best configuration of \toolname depends on the computing platform, we make several observations that can help guide future implementations of \toolname.\\
\textbf{Availability:} \url{https://github.com/CMU-SAFARI/Scrooge}
